\section{Related work}
\label{sec:related_work}

% \subsection{Radiance Fields and Point-based Representations}
% The trajectory of 3D scene reconstruction has evolved from classical geometric formulations, such as Structure-from-Motion (SfM) \cite{schonberger2016structure} and Multi-View Stereo (MVS)~\cite{schonberger2016pixelwise}, toward differentiable neural representations. While early implicit methods like Neural Radiance Fields (NeRF) \cite{mildenhall2021nerf} achieved high-fidelity view synthesis, their reliance on computationally expensive ray-marching limited real-time applicability. This paradigm shifted with the introduction of 3D Gaussian Splatting (3DGS) \cite{kerbl20233d}, which utilizes explicit anisotropic Gaussians to represent the environment. By leveraging tile-based differentiable rasterization, 3DGS facilitates real-time rendering and direct geometric manipulation~\cite{lu2024scaffold,jiang2025horizon,ren2024octree}. 
% To extend this explicit representation to dynamic environments, 4D Gaussian Splatting~\cite{wu20244d} and its variants~\cite{lin2024gaussian,li2024spacetime,zhu2024motiongs} introduce temporal dimensions by modeling Gaussians with time-dependent deformation fields or velocity vectors. These methods enable the high-fidelity reconstruction of non-rigid motions and transient scene elements while maintaining the real-time rendering advantages of the original splatting framework. Furthermore, recent advancements have refined these representations by incorporating 2D Gaussian primitives \cite{huang20242d, dai2024high} to enhance surface coherence and mitigating artifacts through geometric priors such as depth and normal consistency. Such enhancements ensure that these point-based maps are not only visually photorealistic but also structurally accurate for complex downstream spatial and temporal tasks~\cite{jiang2025halogs,ren2025ags,jiang2026planing}.

% \subsection{Deep Visual SLAM and Tracking Foundation Models}
% Simultaneous Localization and Mapping (SLAM) has transitioned from hand-crafted feature tracking to end-to-end differentiable frameworks. DROID-SLAM \cite{teed2021droid} established a benchmark in this domain by employing a Gated Recurrent Unit (GRU) and a differentiable bundle adjustment (BA) layer to iteratively refine camera poses and pixel-wise disparities. 
% Building on learned optimization, MegaSaM~\cite{li2025megasam} addresses the challenges of ``in-the-wild'' videos by integrating monocular depth priors and motion probability networks to robustly handle dynamic objects and low-parallax motion. Moving toward more generalized architectures, MASt3R-SLAM~\cite{murai2024_mast3rslam} leverages dense pointmap regression to enable calibration-free tracking, allowing the framework to operate effectively even with time-varying intrinsics or unknown camera parameters. To further resolve the inherent projective ambiguity in such uncalibrated sequences, VGGT-SLAM~\cite{maggio2025vggt} utilizes factor graph optimization on the $\mathbf{SL}(4)$ manifold, ensuring global consistency across submaps by accounting for 15-DOF projective transformations. These frameworks collectively represent a move toward spatial AI capable of navigating complex, unconstrained environments with minimal prior knowledge.


\subsection{Learning-based and Foundation Model-based SLAM}
In recent years, visual Simultaneous Localization and Mapping (SLAM) has evolved from hand-crafted feature-based pipelines~\cite{campos2021orb,forster2014svo,engel2017direct} toward end-to-end learning-based frameworks~\cite{teed2021droid, teed2023deep}. DROID-SLAM~\cite{teed2021droid} established a strong benchmark by integrating a Gated Recurrent Unit (GRU) with a differentiable bundle adjustment (BA) layer, enabling iterative refinement of camera poses and pixel-wise disparities.
Recent SLAM frameworks increasingly incorporate geometric priors or foundation models to improve robustness and accuracy. For example, MegaSaM~\cite{li2025megasam} addresses the challenges posed by in-the-wild videos by integrating monocular depth priors with motion probability networks, enabling robust handling of dynamic objects and low-parallax motion. Moving toward more generalizable architectures, MASt3R-SLAM~\cite{murai2024_mast3rslam} leverages dense pointmap regression for calibration-free tracking, allowing the framework to operate under unknown camera parameters. To resolve the projective ambiguity inherent in such uncalibrated sequences, VGGT-SLAM~\cite{maggio2025vggt} formulates the optimization problem on the $\mathbf{SL}(4)$ manifold via factor graph optimization, thereby ensuring global consistency across submaps while accounting for 15-DOF projective transformations. %Collectively, these approaches represent an important step toward spatial AI frameworks capable of operating in complex and unconstrained environments with minimal prior knowledge.

\subsection{3D Scene Reconstruction}
The development of 3D scene reconstruction has progressed from classical geometric formulations, such as Structure-from-Motion (SfM) \cite{schonberger2016structure} and Multi-View Stereo (MVS)~\cite{schonberger2016pixelwise}, toward differentiable neural representations. Early implicit approaches, exemplified by Neural Radiance Fields (NeRF)~\cite{mildenhall2021nerf}, demonstrated remarkable fidelity in novel view synthesis; however, their reliance on computationally intensive ray marching limited their applicability in real-time scenarios. This paradigm shifted with the introduction of 3D Gaussian Splatting (3DGS)~\cite{kerbl20233d}, which represents scenes using explicit anisotropic Gaussian primitives. By leveraging tile-based differentiable rasterization, 3DGS enables real-time rendering while supporting explicit optimization~\cite{lu2024scaffold,jiang2025horizon,ren2024octree}.

Recent works have extended 3DGS to support a broader range of tasks. For instance, several studies have refined these representations by introducing 2D Gaussian primitives \cite{huang20242d, dai2024high} to improve surface accuracy, as well as incorporating depth and normal consistency, to mitigate rendering artifacts. Consequently, these point-based representations not only achieve photorealistic rendering quality but also maintain geometric accuracy for complex spatial and temporal tasks.
Moreover, 4D Gaussian Splatting~\cite{wu20244d} and its variants~\cite{lin2024gaussian,li2024spacetime,zhu2024motiongs} incorporate temporal modeling by associating Gaussians with time-dependent deformation fields or velocity vectors. These approaches enable high-fidelity reconstruction of non-rigid motions and transient scene elements while preserving the real-time rendering advantages of the original splatting framework. 

\subsection{Streaming Reconstruction}
Streaming reconstruction aims to incrementally build 3D models from sequential sensor data streams while maintaining low latency~\cite{wang2026towards}. Recent methods extend this paradigm by combining online pose estimation with 3DGS to enable real-time streaming reconstruction. GS-SLAM~\cite{yan2024gs} was among the first to integrate 3DGS into a dense SLAM framework, employing adaptive Gaussian expansion and coarse-to-fine tracking to jointly address mapping and rendering. To enable streaming reconstruction in large-scale environments, Onthefly-NVS~\cite{meuleman2025fly} introduced a pixel spawning mechanism along with a sliding-window anchor strategy, allowing real-time processing of kilometer-scale scenes. For dynamic scenes, Instant4D~\cite{luo2025instant4d} demonstrated minute-level 4D reconstruction from monocular video by employing simplified isotropic Gaussians, reducing redundancy by up to 90\%. 

More recent works further enhance robustness by incorporating priors from geometric or vision foundation models~\cite{artdeco, jiang2026planing, cheng2025outdoor, zhangflash}. 
For example, ARTDECO~\cite{artdeco} combines MASt3R-based feed-forward predictions with a level-of-detail representation to maintain global consistency under unconstrained streaming inputs. 
To address the limited surface coherence of pure Gaussian-based representations, PLANING~\cite{jiang2026planing} employs a hybrid triangle-Gaussian representation within the streaming reconstruction framework. 
By decoupling stable geometric anchors from neural appearance modeling, this design enables structured surface reconstruction suitable for downstream tasks.